\textbf{Valoración global.} La primera respuesta de la IA se renderizó sin errores de sintaxis. Incluía las cuatro relaciones UML exigidas (asociación, agregación, composición y herencia), además de realización y dependencia, las enumeraciones e interfaces solicitadas (\texttt{EstadoPedido}, \texttt{PasarelaPago}, \texttt{ServicioCorreo}) y métodos suficientes para implementar las diez historias. Como ER de partida se le entregó la versión sin corregir, y la IA mantuvo los atributos de todas las entidades. Los problemas detectados son de alcance y coherencia:
\begin{itemize}[leftmargin=*,nosep]
\item \textbf{Atributos inventados, ausentes del ER y del backlog (incumple la regla 4 del prompt).} \texttt{Administradora.codigoEmpleada} no aparece en ninguna historia ni en el ER. \texttt{Cliente.recordarSesion} duplica el parámetro \texttt{recordar} de \texttt{iniciarSesion} y el token de \texttt{Usuario}. \emph{Corrección:} ambos se eliminan.
\item \textbf{Funcionalidad fuera del backlog.} La IA añadió reembolsos (\texttt{reembolsar} en la interfaz y en sus dos implementaciones, y el estado \texttt{REEMBOLSADO}), una notificación de envío por correo (\texttt{enviarNotificacionEnvio}; solo US-4 exige correo, al confirmar la compra) y \texttt{Administradora.actualizarStockTalla} (ninguna historia contempla la gestión de stock). \emph{Corrección:} se eliminan los tres. Se conserva la reserva de stock de \texttt{ProductoTalla}, porque respalda las tallas agotadas de US-2 durante el checkout.
\item \textbf{Responsabilidad duplicada y acoplamiento innecesario.} El envío de un pedido podía marcarse desde \texttt{Administradora.marcarPedidoEnviado(pedido, correo)} y desde \texttt{Pedido.marcarEnviado(correo)}, ambos con \texttt{ServicioCorreo}. Sin la notificación de envío, \emph{corrección:} \texttt{marcarPedidoEnviado(Pedido)} delega en \texttt{Pedido.marcarEnviado()}, sin dependencia del correo.
\item \textbf{Incoherencias derivadas de las correcciones del ER.} \texttt{Usuario.tokenSesion} pasa a \texttt{tokenRecordar} y se añade \texttt{Pedido.pais}, para mantener los nombres alineados con el ER final (regla 4).
\item \textbf{Método de fábrica mal declarado.} \texttt{Pedido.crearDesdeCarrito} se había escrito como método de instancia, pero crea el pedido. \emph{Corrección:} se marca como estático (\texttt{\$}).
\end{itemize}
\textbf{Observaciones aceptadas sin cambios.} (1) La IA usó visibilidad protegida (\texttt{\#}) en las clases abstractas \texttt{Usuario} y \texttt{Pago}, más allá de \texttt{+} y \texttt{-} pedidos en el prompt; se acepta por ser el uso habitual en herencia. (2) Las clases conservan atributos de clave foránea (\texttt{idUsuario}, \texttt{idCupon}, etc.) que duplican asociaciones ya dibujadas; se mantienen por trazabilidad con el ER, aunque en una implementación orientada a objetos se navegaría por la asociación. (3) La clase \texttt{Catalogo} no tiene entidad homóloga en el ER; se acepta como clase de servicio necesaria para US-1 y US-6.
